\documentclass[11pt]{article}

\usepackage[margin=1in]{geometry}
\usepackage{amsmath,amssymb,amsthm,mathtools,mathrsfs}
\usepackage[colorlinks=true,linkcolor=blue,citecolor=blue,urlcolor=blue]{hyperref}

\newcommand{\mH}{\mathcal{H}}
\newcommand{\mA}{\mathcal{A}}
\newcommand{\mB}{\mathcal{B}}
\newcommand{\mL}{\mathcal{L}}
\newcommand{\mM}{\mathcal{M}}
\newcommand{\mK}{\mathcal{K}}
\newcommand{\mN}{\mathcal{N}}
\newcommand{\mE}{\mathcal{E}}
\newcommand{\mR}{\mathcal{R}}
\newcommand{\scrH}{\mathscr{H}}
\newcommand{\scrK}{\mathscr{K}}
\newcommand{\Tr}{\mathrm{Tr}}
\newcommand{\id}{\mathrm{id}}
\newcommand{\barA}{\overline{A}}
\newcommand{\bara}{\overline{a}}
\newcommand{\ket}[1]{|#1\rangle}
\newcommand{\bra}[1]{\langle #1|}
\newcommand{\braket}[2]{\langle #1|#2\rangle}
\newcommand{\code}{\mathrm{code}}
\newcommand{\wt}[1]{\widetilde{#1}}
\newcommand{\nn}{\nonumber}

\newtheorem{theorem}{Theorem}
\newtheorem{definition}{Definition}
\newtheorem{remark}{Remark}

\title{Section 2: QEC and Algebra Inclusions}
\author{}
\date{}

\begin{document}
\maketitle

\subsection{Subspace QECC}

Consider a finite-dimensional quantum system represented in Hilbert space $\mH$ with the algebra of observables $\mA$ and denote by $\mathcal{L}_1(\mH)$ the set of trace class operators on $\mH$. In the Schrodinger picture, a quantum channel is a trace-preserving CP map $\Phi^*:\mathcal{L}_1(\mH)\to \mathcal{L}_1(\mH)$. It admits a Kraus representation:
\begin{align}
    &\forall\rho\in \mathcal{L}_1(\mH):\qquad \Phi^*(\rho)=\sum_i K_i \rho K_i^\dagger\nn\\
    &\sum_i K_i K_i^\dagger=1
\end{align}


In standard QEC, we consider a {\it code subspace} $\mH_C\subset \mH$, and we say this code is correctible with respect to the quantum channel $\Phi^*:\mathcal{L}_1(\mH)\to \mathcal{L}_1(\mH)$ 

\subsection{Operator algebraic QECC}

Consider the observable algebra of a {\it finite} quantum system $\mA$ acting on the Hilbert space $\mK$. The center $Z(\mA)$ is an Abelian algebra generated by a unique set of mutually orthogonal projections $e_r$, i.e., $e_r e_{r'}=\delta_{rr'}e_r$ with $\sum_r e_r=1$. This means that we can decompose $\mA=\oplus_r \mA_r$ with $\mA_r=e_r\mA=\mA e_r=e_r \mA e_r$. 

Consider a general noise map in the Heisenberg picture $\Phi:\mA\to \mA$. 
It is a unital CP map, which means that it admits a Stinespring dilation:
\begin{align}
    \Phi(a)=V^\dagger \pi(a) V
\end{align}
for isometries $V:\mK\to \mH$, i.e., $V^\dagger V=1$, and $\pi$ a $*$-homomorphism, i.e., a representation of $\mA$ on $B(\mH)$. The projection $VV^\dagger=P_C$ projects to a subspace we will call {\it code-subspace}:
\begin{align}
    \mH_C\equiv P_C\mH=VV^\dagger\mH=V\mK\ .
\end{align}
In finite dimensions, we have $Z(\pi(\mA))=\pi(Z(\mA))$. 

\paragraph{Trivial center:} For simplicity, assume that $\mA$ is a factor so that its center is trivial. Then, the Stinespring dilatation takes the form
\begin{align}
 %&   \Phi(a)=V^\dagger U^\dagger(a\otimes 1_{\bar{A}}) U V\nn\\
 &\pi(a)=U_{A\bar{A}}(a\otimes 1_{\bar{A}}) U_{A\bar{A}},\qquad V=1_A\otimes \ket{0}_{\bar{A}}, \qquad P_C=1_A\otimes \ket{0}\bra{0}_{\bar{A}}\nn\\
 &\Phi(a)=\text{tr}_{\bar{A}}(P_C U^\dagger (a\otimes 1) U P_C)=W^\dagger(a\otimes 1) W, \qquad W=U V:\mH_K\to \mH_A\otimes \mH_{\bar{A}}
\end{align}
where $\mH=\mH_A\otimes \mH_{\bar{A}}$ with $\bar{A}$ some fiducial environment, $\mH_A\simeq \mK$, and $V$ is an isometry that we will interpret as {\it encoding} in the language of Quantum Error Correction Codes (QECC) with 
\begin{align}
    &\ket{\psi_i}=V\ket{i}_K=\ket{i}_A\otimes \ket{0}_{\bar{A}}
\end{align}
Consider the subalgebra of operators $\mA^\Phi\in \mA$ such that
\begin{align}
    \forall c\in \mA^C: \qquad \pi(c)V=V c
\end{align}
The action of the noise map $\Phi$ restricted to $\mA^\Phi$ is multiplicative (endormophism):
\begin{align}
    \Phi(c_1)\Phi(c_2)=V^\dagger \pi(c_1) P_C \pi(c_2) V=V^\dagger P_C \pi(c_1)\pi(c_2)V=V^\dagger \pi(c_1c_2)V=\Phi(c_1c_2)\ .
\end{align}
This subalgebra is called the {\it multiplicative} domain of $\Phi$ in $\mA$. In the language of QECC, $\mA^\Phi$ is the unitarily correctible subsystem. 

*Add the definition of a correctible subsystem and unitarily correctible subsystem.

Then, generalize to when $\mA$ has a center, and we obtain the general form of a QECC as in Harlow's paper, for erasure noise.

The action of errors on code states $\ket{\psi_i}$ can be unitarily undone by $U^\dagger$. The algebra $\mA$ is the unitarily correctible algebra of the noise map $\Phi$. 



A conceptual advantage of Stinespring dilation is that it realizes an arbitrary noise map as erasure, or more precisely, the conditional expectation that is the restriction of $\mE:B(\mH)\to \pi(\mA)$. To understand this, start with the case where $\mA$ is a factor. Then, the inclusion of factors $\pi(\mA)\subset B(\mH)$ implies that there exists a decomposition 
\begin{align}
    \mH=(\mH_A\otimes \mH_{A'})\oplus \mH_E
\end{align}
such that $\pi(\mA)=B(\mH_A)$. In other words, there exists a unitary $U:\mH\to \mH$ such that 
\begin{align}
\Phi(a)=V^\dagger U^\dagger (a\otimes 1)U V=\text{tr}_{A'}(U^\dagger (a\otimes 1) U)
\end{align}
The noise map is the erasure of $\bar{A}$:

and the set of all operators $\pi(\mA)=B(\mH_A)$ forms a correctible (decoherence-free) subalgebra, because every operator in $\pi(a)$ commutes with 

In matrix algebras, any inclusion of algebras $\pi(\mA)\subset B(\mH)$ implies that a decomposition:
\begin{align}
    &\mH=\oplus_r (\mH_r\otimes \mH'_r)\oplus \mH_E\nn\\
    &\pi(\mA)=\oplus_r \pi_r(\mA)\otimes 1'_{r}
\end{align}
For matrix algebras, we have
\begin{align}
    Z(\pi(\mA))=\pi(Z(\mA))
\end{align}

For instance, in the case of matrix algebras, if $\pi(\mA)$ is a factor:
\begin{align}
    \Phi(a)=V^\dagger (a\otimes 1) V\ .
\end{align}
which means that the code subspace decomposes as $\mH_C=\mH_A\otimes \mH_{\bar{A}}$ such that $\pi(\mA)\subseteq B(\mH_A)$, and the noise map is simply erasure of $\bar{A}$. In the general case, $\pi(\mA)$ has a non-trivial center, we have the decomposition:
\begin{align}
&\mH_C=\oplus_r \mH_{A,r}\otimes \mH_{\bar{A},r}\nn\\
&\pi(\mA)=\oplus_r \pi_r(\mA_A)\otimes 1_{\bar{A},r}
\end{align}
Viewing noise as erasure makes the concept of {\it correctible subalgebra} explicit. In a general QECC, we perform syndrome measurements and based on the outcome of the measurement, we perform different recovery maps. A correctible subalgebra is a subalgebra that can be corrected with a single recovery map (no syndrom measurement). More explicitly, we have the definition:
\begin{definition}
A $C^*$-subalgebra $\mA^C\subset \mA$ is called {\it correctible} if there exists a unital CP recovery map $\mathcal{R}$:
\begin{align}
    \forall a\in \mA^C:\qquad \mR\circ \Phi(a)=a
\end{align}
\end{definition}
Consider an orthonormal basis for 


\section*{Summary of Harlow Section 5 up to Section 5.2}

This note summarizes the part of Daniel Harlow's
\emph{The Ryu--Takayanagi Formula from Quantum Error Correction}
\cite{HarlowRTQEC} beginning with Section 5 and ending just before Section 5.2.
The purpose of that part of the paper is to pass from subsystem quantum error
correction to \emph{operator-algebra quantum error correction}: instead of
requiring the recovery of all logical operators, one fixes a finite-dimensional
von Neumann algebra \(M\subset \mB(\mH_{\code})\) and asks only for the recovery
of the operators in \(M\).  This is the natural setting for entanglement wedge
reconstruction when the relevant bulk algebra has a nontrivial center.

\subsection*{Finite-dimensional von Neumann algebras}

Let \(\mH_{\code}\) be finite-dimensional and let \(M\subset \mB(\mH_{\code})\)
be a unital \(*\)-subalgebra.  In finite dimension this is a von Neumann
algebra.  The finite-dimensional classification theorem gives a canonical block
decomposition
\begin{equation}
    \mH_{\code}
    =
    \bigoplus_{\alpha}
    \left(\mH_{a_\alpha}\otimes \mH_{\bara_\alpha}\right),
    \label{eq:code-block-decomp}
\end{equation}
such that
\begin{align}
    M
    &=
    \bigoplus_{\alpha}
    \left(\mB(\mH_{a_\alpha})\otimes I_{\bara_\alpha}\right),
    \label{eq:M-block-form}\\
    M'
    &=
    \bigoplus_{\alpha}
    \left(I_{a_\alpha}\otimes \mB(\mH_{\bara_\alpha})\right),
    \label{eq:commutant-block-form}\\
    Z(M)
    &=
    M\cap M'
    =
    \left\{
    \bigoplus_\alpha
    \lambda_\alpha I_{a_\alpha}\otimes I_{\bara_\alpha}
    :
    \lambda_\alpha\in\mathbb C
    \right\}.
    \label{eq:center-block-form}
\end{align}
Thus each block \(\alpha\) is a superselection sector associated with the
center.  Inside a fixed \(\alpha\)-block, \(M\) acts only on the
\(a_\alpha\)-factor, while \(M'\) acts only on the \(\bara_\alpha\)-factor.
The special case in which there is only one block is precisely a factor
algebra, recovering ordinary subsystem QEC.

For a state \(\wt{\rho}\) on \(\mH_{\code}\), write
\(\wt{\rho}_{\alpha\alpha}\) for its diagonal block in
\eqref{eq:code-block-decomp}.  Define
\begin{align}
    p_\alpha \wt{\rho}_{a_\alpha}
    &:=
    \Tr_{\bara_\alpha}\wt{\rho}_{\alpha\alpha},
    &
    p_\alpha \wt{\rho}_{\bara_\alpha}
    &:=
    \Tr_{a_\alpha}\wt{\rho}_{\alpha\alpha},
    \label{eq:block-reduced-states}
\end{align}
with \(p_\alpha\ge 0\), \(\sum_\alpha p_\alpha=1\), and each reduced density
matrix normalized when \(p_\alpha\ne0\).  The algebraic entropies are
\begin{align}
    S(\wt{\rho},M)
    &=
    -\sum_\alpha p_\alpha\log p_\alpha
    +
    \sum_\alpha p_\alpha S(\wt{\rho}_{a_\alpha}),
    \label{eq:alg-entropy-M}\\
    S(\wt{\rho},M')
    &=
    -\sum_\alpha p_\alpha\log p_\alpha
    +
    \sum_\alpha p_\alpha S(\wt{\rho}_{\bara_\alpha}).
    \label{eq:alg-entropy-Mprime}
\end{align}
The first term is the Shannon entropy of the central sector, and the second is
the sector-averaged ordinary von Neumann entropy.

\subsection*{A rigorous form of Harlow's Theorem 5.1}

Let the physical Hilbert space factorize as
\[
    \mH=\mH_A\otimes\mH_{\barA},
\]
and let \(\mH_{\code}\subset \mH\) be a code subspace with projection
\(P_{\code}\).  Let \(V:\mH_{\code}\hookrightarrow \mH_A\otimes\mH_{\barA}\)
denote the inclusion isometry.  Choose orthonormal bases
\(\{\ket{\alpha,i}_{a_\alpha}\}\) and
\(\{\ket{\alpha,j}_{\bara_\alpha}\}\), and write
\[
    \ket{\wt{\alpha,ij}}
    :=
    V\left(
    \ket{\alpha,i}_{a_\alpha}\otimes
    \ket{\alpha,j}_{\bara_\alpha}
    \right).
\]
Let \(R\) be a reference system isomorphic to \(\mH_{\code}\), with basis
\(\{\ket{\alpha,ij}_R\}\), and define the maximally entangled state between
the reference and the code:
\begin{equation}
    \ket{\Phi}
    =
    \frac{1}{\sqrt{|\mH_{\code}|}}
    \sum_{\alpha,i,j}
    \ket{\alpha,ij}_R
    \ket{\wt{\alpha,ij}}_{A\barA}.
    \label{eq:reference-state}
\end{equation}
For \(O\in M\), define \(O_R\) to be the operator on \(R\) with the same matrix
elements as \(O^T\) in the chosen basis, equivalently the unique operator
satisfying
\[
    O_R\ket{\Phi}=O\ket{\Phi},
    \qquad
    O_R^\dagger\ket{\Phi}=O^\dagger\ket{\Phi}.
\]

\begin{theorem}[Operator-algebra erasure correction]
\label{thm:harlow-oa-qec}
With the setup above, the following four statements are equivalent.

\begin{enumerate}
    \item \textbf{Structural decoding form.}
    There is a decomposition
    \begin{equation}
        \mH_A
        =
        \bigoplus_\alpha
        \left(
        \mH_{A_1^\alpha}\otimes \mH_{A_2^\alpha}
        \right)
        \oplus \mH_{A_3},
        \qquad
        \dim \mH_{A_1^\alpha}=\dim \mH_{a_\alpha},
        \label{eq:A-decomp}
    \end{equation}
    so in particular \(\sum_\alpha \dim \mH_{a_\alpha}\leq \dim \mH_A\), a
    unitary \(U_A\) on \(\mH_A\), and for every \(\alpha\) a set of orthonormal
    states
    \[
        \ket{\chi_{\alpha,j}}_{A_2^\alpha\barA}
        \in
        \mH_{A_2^\alpha}\otimes \mH_{\barA}
    \]
    such that
    \begin{equation}
        \ket{\wt{\alpha,ij}}
        =
        U_A
        \left(
        \ket{\alpha,i}_{A_1^\alpha}
        \otimes
        \ket{\chi_{\alpha,j}}_{A_2^\alpha\barA}
        \right).
        \label{eq:structure-decoding}
    \end{equation}

    \item \textbf{Logical reconstruction on \(A\).}
    For every \(O\in M\), there exists an operator \(O_A\in\mB(\mH_A)\) such
    that for every \(\ket{\wt{\psi}}\in\mH_{\code}\),
    \begin{equation}
        O_A\ket{\wt{\psi}}=O\ket{\wt{\psi}},
        \qquad
        O_A^\dagger\ket{\wt{\psi}}=O^\dagger\ket{\wt{\psi}}.
        \label{eq:reconstruction-A}
    \end{equation}

    \item \textbf{Environment algebra compresses into the commutant.}
    For every \(X_{\barA}\in\mB(\mH_{\barA})\), the compressed operator
    \(P_{\code}X_{\barA}P_{\code}\), viewed as an operator on
    \(\mH_{\code}\), lies in \(M'\):
    \begin{equation}
        P_{\code}X_{\barA}P_{\code}
        =
        X'P_{\code}
        \qquad
        \text{for some } X'\in M'.
        \label{eq:compressed-environment}
    \end{equation}

    \item \textbf{Reference--environment commutation criterion.}
    For every \(O\in M\),
    \begin{equation}
        [O_R,\rho_{R\barA}(\Phi)]=0,
        \label{eq:reference-criterion}
    \end{equation}
    where \(\rho_{R\barA}(\Phi)=\Tr_A\ket{\Phi}\bra{\Phi}\).
\end{enumerate}
\end{theorem}

\paragraph{Interpretation.}
Condition \eqref{eq:reconstruction-A} says that the logical algebra \(M\) is
recoverable after erasing \(\barA\).  Condition
\eqref{eq:compressed-environment} is the Heisenberg-picture no-disturbance
criterion: operators on the erased subsystem can act nontrivially on the code,
but only through the commutant \(M'\), never through the protected algebra
\(M\).  Condition \eqref{eq:reference-criterion} is the information-disturbance
formulation: the joint state of the reference and erased subsystem may retain
information, but it must commute with the reference copy of \(M\).  Thus the
erased subsystem can know about \(M'\) and about the central label \(\alpha\),
but not about the quantum degrees of freedom protected by \(M\).

\paragraph{Proof skeleton.}
The implication \((1)\Rightarrow(2)\) is constructive.  If \(O\in M\) has block
form \(O=\oplus_\alpha (O_{a_\alpha}\otimes I_{\bara_\alpha})\), define
\[
    O_A
    =
    U_A
    \left[
    \bigoplus_\alpha
    \left(
    O_{A_1^\alpha}\otimes I_{A_2^\alpha}
    \right)
    \right]
    U_A^\dagger,
\]
where \(O_{A_1^\alpha}\) has the same matrix elements as
\(O_{a_\alpha}\).  Equation \eqref{eq:structure-decoding} then gives
\eqref{eq:reconstruction-A}.

For \((2)\Rightarrow(3)\), suppose the compressed environment operator
\(x'=P_{\code}X_{\barA}P_{\code}\) were not in \(M'\).  Then there would be
some \(O\in M\) with \([x',O]\ne0\).  Hence there are code vectors
\(\ket{\wt{\psi}}\) and \(\ket{\wt{\eta}}\) with
\[
    \bra{\wt{\psi}}[x',O]\ket{\wt{\eta}}\ne0.
\]
But \(X_{\barA}\) commutes with every reconstruction \(O_A\) on \(A\), and
condition (2) identifies the matrix elements of \(O_A\) and \(O\) on the code.
This forces all such code matrix elements of \([x',O]\) to vanish, a
contradiction.  Hence \(x'\in M'\).

For \((3)\Rightarrow(4)\), insert arbitrary test operators
\(X_{\barA}\) and \(Y_R\) and use \eqref{eq:compressed-environment}:
\[
    \Tr_{R\barA}
    \left(
    O_R\rho_{R\barA}(\Phi)X_{\barA}Y_R
    \right)
    =
    \Tr_{R\barA}
    \left(
    \rho_{R\barA}(\Phi)O_RX_{\barA}Y_R
    \right).
\]
Since this holds for all \(X_{\barA}\) and \(Y_R\), the density matrix
\(\rho_{R\barA}(\Phi)\) commutes with \(O_R\).

For \((4)\Rightarrow(1)\), decompose the reference as
\[
    \mH_R
    =
    \bigoplus_\alpha
    \left(
    \mH_{R_\alpha}\otimes\mH_{\bar R_\alpha}
    \right),
    \qquad
    \dim R_\alpha=\dim a_\alpha,\quad
    \dim \bar R_\alpha=\dim \bara_\alpha .
\]
Because \(O_R\) ranges over
\(\oplus_\alpha(\mB(R_\alpha)\otimes I_{\bar R_\alpha})\), the commutation
condition \eqref{eq:reference-criterion} forces
\begin{equation}
    \rho_{R\barA}(\Phi)
    =
    \bigoplus_\alpha
    \frac{|R_\alpha||\bar R_\alpha|}{|\mH_{\code}|}
    \left(
    \frac{I_{R_\alpha}}{|R_\alpha|}
    \otimes
    \rho_{\bar R_\alpha\barA}^{(\alpha)}
    \right).
    \label{eq:rho-block-form}
\end{equation}
The marginal on \(R\) is maximally mixed, so
\[
    \Tr_{\barA}\rho_{\bar R_\alpha\barA}^{(\alpha)}
    =
    \frac{I_{\bar R_\alpha}}{|\bar R_\alpha|}.
\]
Since \(\ket{\Phi}\) purifies \(\rho_{R\barA}(\Phi)\) on \(A\), the Schmidt
rank condition implies
\[
    \sum_\alpha
    |R_\alpha|\,
    \operatorname{rank}\rho_{\bar R_\alpha\barA}^{(\alpha)}
    \le
    |A|.
\]
One can therefore choose the direct-sum decomposition
\eqref{eq:A-decomp} with enough room in each \(A_2^\alpha\) to purify
\(\rho_{\bar R_\alpha\barA}^{(\alpha)}\).  Such a purification can be written
as
\[
    \ket{\psi_\alpha}_{\bar R_\alpha A_2^\alpha\barA}
    =
    \frac{1}{\sqrt{|\bar R_\alpha|}}
    \sum_j
    \ket{\alpha,j}_{\bar R_\alpha}
    \ket{\chi_{\alpha,j}}_{A_2^\alpha\barA},
\]
and the maximally mixed marginal on \(\bar R_\alpha\) implies that the
\(\ket{\chi_{\alpha,j}}\) are orthonormal.  Combining these purifications gives
a state of the form \eqref{eq:structure-decoding}; since it and
\(\ket{\Phi}\) purify the same \(\rho_{R\barA}(\Phi)\), they differ by a
unitary on \(A\).  This gives condition (1).

\subsection*{Complementary recovery and algebra inclusions}

For holographic applications Harlow imposes a complementary recovery condition:
not only should \(M\) be reconstructable on \(A\), but \(M'\) should be
reconstructable on \(\barA\).  Applying Theorem
\ref{thm:harlow-oa-qec} to both sides gives the normal form
\begin{equation}
    \ket{\wt{\alpha,ij}}
    =
    U_AU_{\barA}
    \left(
    \ket{\alpha,i}_{A_1^\alpha}
    \ket{\alpha,j}_{\barA_1^\alpha}
    \ket{\chi_\alpha}_{A_2^\alpha\barA_2^\alpha}
    \right),
    \label{eq:complementary-normal-form}
\end{equation}
with
\[
    \dim A_1^\alpha=\dim a_\alpha,
    \qquad
    \dim \barA_1^\alpha=\dim \bara_\alpha.
\]
This is the operator-algebraic version of the subsystem-code circuit picture.
After local decoding unitaries, the logical algebra inclusion is visible block
by block:
\[
    M
    \hookrightarrow
    \bigoplus_\alpha
    \mB(\mH_{A_1^\alpha})\otimes I_{A_2^\alpha},
    \qquad
    M'
    \hookrightarrow
    \bigoplus_\alpha
    I_{\barA_1^\alpha}\otimes \mB(\mH_{\barA_1^\alpha}).
\]
The center is the shared classical label \(\alpha\).  The entangled states
\(\ket{\chi_\alpha}_{A_2^\alpha\barA_2^\alpha}\) are not logical degrees of
freedom; rather, they are the universal entanglement resource that later
produces the area term in Section 5.2.

\subsection*{Faulkner Section 2: the same theorem as an algebra inclusion}

Faulkner's Section 2 of \emph{The holographic map as a conditional expectation}
\cite{FaulknerConditional} abstracts the preceding finite-dimensional theorem
away from the explicit block decomposition
\eqref{eq:code-block-decomp}--\eqref{eq:center-block-form}.  The code is an
isometry
\[
    V:\scrH=\scrH_{\code}\longrightarrow \scrK=\scrH_{\mathrm{phys}},
    \qquad
    V^\dagger V=1_{\scrH}.
\]
The code algebra is a von Neumann algebra
\(\mN\subset \mB(\scrH)\), and the physical algebra from which one wants to
reconstruct is a von Neumann algebra \(\mM\subset\mB(\scrK)\).  The notation is
related to the Harlow notation above by the dictionary
\[
    \mN \leftrightarrow M,
    \qquad
    \mN' \leftrightarrow M',
    \qquad
    \mM \leftrightarrow \mB(\mH_A)\otimes I_{\barA},
    \qquad
    \mM' \leftrightarrow I_A\otimes \mB(\mH_{\barA})
\]
in the simplest finite-dimensional erasure problem.  More generally
\(\mM\) is the chosen boundary algebra for a region and \(\mM'\) is its
commutant, interpreted as the algebra of the complementary environment.

For a vector \(\psi\), Faulkner writes
\(\omega_\psi(x)=\langle\psi,x\psi\rangle\) for the associated normal vector
functional.  His operator-algebra error-correction theorem may then be stated
as follows.

\begin{theorem}[Faulkner's reconstruction theorem]
\label{thm:faulkner-oa-qec}
Let \(V:\scrH\to\scrK\) be an isometry, let
\(\mN\subset\mB(\scrH)\) and \(\mM\subset\mB(\scrK)\) be
\(\sigma\)-finite von Neumann algebras, and use commutants relative to their
given representations.  Up to the standard support-projection qualifications
in the nonunital case, the following conditions are equivalent.

\begin{enumerate}
    \item \textbf{DHW privacy condition.}  For all vectors
    \(\psi,\phi\in\scrH\),
    \begin{equation}
        \omega_\psi|_{\mN'}=\omega_\phi|_{\mN'}
        \quad\Longrightarrow\quad
        \omega_{V\psi}|_{\mM'}=\omega_{V\phi}|_{\mM'} .
        \label{eq:faulkner-dhw}
    \end{equation}

    \item \textbf{Compressed environment channel.}  The Heisenberg compression
    of the physical commutant lands in the code commutant:
    \begin{equation}
        \alpha'(m'):=V^\dagger m'V\in\mN',
        \qquad
        m'\in\mM',
        \label{eq:faulkner-alpha-prime}
    \end{equation}
    so \(\alpha':\mM'\to\mN'\) is a normal unital completely positive map.

    \item \textbf{Reconstruction embedding.}  There exists a normal injective
    \(*\)-homomorphism
    \begin{equation}
        \beta:\mN\longrightarrow \mM
        \label{eq:faulkner-beta}
    \end{equation}
    such that every logical operator \(n\in\mN\) is represented on the
    physical algebra by
    \begin{equation}
        \beta(n)V=Vn.
        \label{eq:faulkner-intertwining}
    \end{equation}
\end{enumerate}
Faulkner gives both a nonunital support-projection version and a standard
unital version.  In the standard version one assumes, for example, the
existence of a code vector \(\eta\) for which \(V\eta\) is cyclic and
separating for \(\mM\); then \(\alpha'\) is faithful, \(\beta\) is unital, and
\(\beta\) is unique.
\end{theorem}

This is the von Neumann algebra version of Theorem
\ref{thm:harlow-oa-qec}.  In the finite-dimensional setting, the map
\(\alpha'\) in \eqref{eq:faulkner-alpha-prime} is precisely the compression in
\eqref{eq:compressed-environment}: if \(m'=X_{\barA}\) acts only on the erased
subsystem, then
\[
    V^\dagger X_{\barA}V\in M'
\]
is the statement that the environment can see only the commutant algebra.
Likewise, the embedding condition \eqref{eq:faulkner-intertwining} is the
coordinate-free form of reconstruction on \(A\):
\[
    O_A\ket{\wt{\psi}}=O\ket{\wt{\psi}}
    \quad\text{becomes}\quad
    \beta(n)V=Vn .
\]
Finally, the reference criterion
\eqref{eq:reference-criterion} is Harlow's finite-dimensional way of packaging
the DHW implication \eqref{eq:faulkner-dhw}: states that agree on \(M'\) may
differ in the protected \(M\)-degrees of freedom, but such differences are not
visible to the physical commutant \(\mM'\).

The important conceptual shift is that Faulkner does not use the direct-sum
normal form of finite-dimensional algebras as the proof engine.  Instead, the
theorem is formulated in the intrinsic language of normal completely positive
maps and normal \(*\)-homomorphisms.  This makes the statement suitable for
infinite-dimensional and type-III settings, where a decomposition into
\(\mH_{a_\alpha}\otimes\mH_{\bara_\alpha}\) sectors is generally unavailable.
It also cleanly separates two maps that are sometimes conflated.  The
compression
\[
    \alpha(m):=V^\dagger mV,\qquad m\in\mM,
\]
is the noisy Heisenberg channel from the physical algebra back to the code
Hilbert space, while \(\beta\) is the reconstruction embedding from the code
algebra into the physical algebra.  Section 2 proves the existence of
\(\beta\) from the privacy condition; it does not yet require a Petz-map
argument or equality of relative entropies.

This distinction is also where conditional expectations enter the later
paper.  Section 2 fixes notation for conditional expectations, but the theorem
itself only gives the embedding \(\beta:\mN\to\mM\) and the compression maps
\(\alpha,\alpha'\).  A genuine conditional expectation of the form
\[
    E=\beta\circ\alpha:\mM\longrightarrow \beta(\mN)
\]
appears once one imposes complementary recovery and applies the same theorem
to the commutant algebra.  Thus, relative to the Harlow discussion above,
Faulkner's Section 2 is best understood as the rigorous algebra-inclusion
version of exact entanglement-wedge reconstruction, while the subsequent
conditional-expectation structure is the extra consequence of complementary
recovery.

\subsection*{Minimal example}

The example at the end of Harlow's Section 5.1 is the two-qubit repetition
code
\[
    \ket{\wt{0}}=\ket{00},
    \qquad
    \ket{\wt{1}}=\ket{11}.
\]
Let \(M\) be the abelian algebra generated by the logical identity and logical
\(\wt{Z}\), where
\[
    \wt{Z}\ket{\wt{0}}=\ket{\wt{0}},
    \qquad
    \wt{Z}\ket{\wt{1}}=-\ket{\wt{1}}.
\]
Here \(M=M'=Z(M)\): the algebra is purely central and has no noncommutative
matrix block.  The logical \(\wt{Z}\) is reconstructable on either physical
qubit, since \(Z_1\) and \(Z_2\) both act as \(\wt{Z}\) on the code subspace.
This example illustrates why algebraic QEC is strictly more flexible than
factor/subsystem QEC: a central observable can be redundantly accessible on both
sides without either side recovering the full logical state.

\begin{thebibliography}{9}
\bibitem{HarlowRTQEC}
D. Harlow,
\emph{The Ryu--Takayanagi Formula from Quantum Error Correction},
arXiv:1607.03901.
\url{https://arxiv.org/abs/1607.03901}.

\bibitem{FaulknerConditional}
T. Faulkner,
\emph{The holographic map as a conditional expectation},
arXiv:2008.04810.
\url{https://arxiv.org/abs/2008.04810}.
\end{thebibliography}

\end{document}
